\subsection{Zero-knowledge circuit construction}
To certify a four-valued decision without revealing the original observations, the circuit checks signatures, coordinate bounds, and context consistency while determining whether threshold-supported states exist on either side of the policy boundary. We implement two geometric backends: lower-corner maximum-depth reachability and threshold-subset enumeration. Their signature-verification and public-output constraints are shared; they differ in how threshold reachability is calculated.

\subsubsection{Lower-corner maximum-depth reachability}
For a query box $Q\subseteq\Omega$, first clip every evidence box $U_i$ to $Q$. For each axis $a$, let $A_a$ contain the lower endpoints of the nonempty clipped boxes. The following lemma reduces maximum-depth search over the entire query region to a finite candidate set.

\begin{lemma}[Lower-corner maximum-depth reachability]\label{lem:corner}
If the maximum overlap depth inside $Q$ is positive, then some point attaining that depth belongs to $A_1\times\cdots\times A_d$.
\end{lemma}
\begin{proof}
Choose a point $z\in Q$ with maximum depth and let $\mathcal S$ be the collection of clipped boxes containing $z$. Their intersection is nonempty. For each axis $a$, define $v_a$ as the greatest lower endpoint of the boxes in $\mathcal S$. This endpoint belongs to $A_a$. Because $z$ belongs to every box in $\mathcal S$, the greatest lower endpoint on any axis cannot exceed the smallest corresponding upper endpoint. Hence $v=(v_1,\ldots,v_d)$ lies in every box in $\mathcal S$ and has depth at least $|\mathcal S|$, the maximum depth attained by $z$. By maximality, its depth is exactly the same. Thus $v$ is a maximum-depth point in the stated candidate set.
\end{proof}

By the lemma, the circuit enumerates at most $n^d$ lower-corner combinations, including duplicates, and checks both membership in $Q$ and coverage by the signed evidence boxes. A point with coverage at least $k$ exists in $Q$ if and only if an enumerated candidate meets that threshold. A candidate contributed by an interval not intersecting $Q$ cannot automatically be treated as valid: the necessary membership tests still apply.

When $P$ is an axis-aligned integer box, $\Omega\setminus P$ is the union of at most $2d$ axial slabs. Along axis $a$, a lower slab ends at $p_a^- -1$ and an upper slab starts at $p_a^+ +1$; all other axes span the discrete domain. Slabs lying outside $\Omega$ are inactive. The circuit calculates $I$ from reachability in $P$ and $O$ from the disjunction of reachability over complement slabs. It then constrains the public decision code to
\begin{equation}\label{eq:decisioncode}
\delta=(1-I)O+2IO+3(1-I)(1-O).
\end{equation}
For Boolean $I$ and $O$, this expression is exactly the coding in Table~\ref{tab:decisions}. Further constraints check the coordinate bit width, ordering of lower and upper endpoints, and containment of boxes and policy within $\Omega$. EdDSA/MiMC signatures are verified inside the circuit, and issuer public keys are compared for pairwise distinctness using their actual encoded elliptic-curve representations, rather than treating a public key as a single scalar.

\subsubsection{Threshold-subset enumeration and computational cost}
An alternative representation of the robust feasible set is

\begin{equation}\label{eq:subset}
\mathcal{F}_k = \bigcup_{\substack{S \subseteq [n] \\ |S| = k}} \left( \bigcap_{i \in S} U_i \right), \quad [n] = \{1, \dots, n\}.
\end{equation}


The intersection of $k$ axis-aligned integer boxes is either another box or the empty set. Its lower and upper endpoints are obtained via coordinatewise maxima and minima, enabling direct tests for nonemptiness, intersection with $P$, and reachability outside $P$. Both the native reference classifier and the second circuit backend use this exact identity, requiring $O(\binom{n}{k}kd)$ geometric operations.

For a fixed dimension $d$, lower-corner reachability requires $O((2d + 1)dn^{d+1})$ geometric comparisons alongside $n$ signature verifications and $O(n^2)$ pairwise key-distinctness checks, and accepts any admissible threshold $k$ as a dynamic public input without recompilation. In contrast, threshold-subset enumeration fixes $k$ at compile time and constrains it against the public threshold. When the total omission budget $b = n - k = f + r$ is a small constant, its geometric work scales as $O(n^{b+1}d)$, whereas as $k$ approaches $n/2$, the subset count $\binom{n}{k}$ grows combinatorially.

For example, at $(n, d, k) = (5, 3, 4)$, subset enumeration checks $\binom{5}{4} = 5$ intersections, whereas lower-corner reachability evaluates up to $5^3 = 125$ candidates per query region. Using the geometric operation proxies $W_D = (2d + 1)dn^{d+1}$ and $W_S = \binom{n}{k}kd$, we have $W_D = 3,360,000$ versus $W_S = 1,140$ at $(n, d, k) = (20, 3, 19)$; when $k$ decreases to $10$, $W_D$ remains constant while $W_S$ increases to $5,542,680$. Thus, both backends implement the exact same four-valued relation while trading runtime threshold flexibility against geometric circuit size in small-omission regimes.
